\begin{textsample}{POS Dim 1 – 2017\_c5 – Score 16.00 – 2017\_c5/enade\_2017\_eng\_amb\_82fdb2c61c.txt}  \label{ex:f1_pos_054}
CAPÍTULO 5 DISTRIBUIÇÃO DOS CONCEITOS 5.1 \textbf{PANORAMA} NACIONAL DA DISTRIBUIÇÃO DOS CONCEITOS O Conceito Enade é um indicador de qualidade que avalia os cursos por \textbf{intermédio} dos desempenhos dos estudantes a partir dos resultados do Enade.
Os valores possíveis vão de 1 ( pior situação ) a 5 ( melhor situação ).
Os cursos com apenas um ou sem nenhum concluinte participante não obtêm o Conceito Enade, ficando Sem Conceito ( SC ).
A Tabela 5.1 apresenta a quantidade e distribuição de cursos de Engenharia Ambiental participantes do Enade/2017, por faixa de conceito e Grande Região.
Enfatiza -se, mais uma vez, que a diferença entre os cursos \textbf{tabulados} neste capítulo e no capítulo 2 corresponde aos cursos Sem Conceito ( SC ), ou seja àqueles, em princípio, sem alunos concluintes que houvessem participado da prova ou àqueles que tivessem tão somente um aluno concluinte18.
Observando -se os dados da Tabela 5.1, nota -se que, dos 252 cursos participantes, 104 ( 41,3 \% ) classificaram -se com conceito 3, o valor modal.
Esse foi também o conceito modal em todas as regiões : Norte ( 50,0 \% ), Nordeste ( 50,0 \% ), Sudeste ( 37,5 \% ), Sul ( 42,3 \% ) e Centro-Oeste ( 40,0 \% ).
O conceito 4 foi o segundo mais frequente em nível nacional ( 22,2 \%, correspondendo a 56 cursos ) e o conceito 2, o terceiro ( 21,4 \%, correspondendo a 54 cursos ).
Houve, ainda, 21 cursos ( 8,3 \% ) que receberam conceito 5 e 12 cursos ( 4,8 \% ) que receberam conceito 1.
Dos 252 cursos de Engenharia Ambiental, cinco ( 2,0 \% ) ficaram Sem Conceito ( SC ).
A região Norte participou com 16 cursos ou 6,3 \% do total nacional.
Desses, 8 cursos receberam o conceito 3, o que \textbf{equivale} a 50,0 \% do total regional, além de ser o conceito 18 Estes cursos com somente um concluinte não foram considerados no capítulo 2. modal na região, como já comentado.
A três cursos ( 18,8 \% ) atribuiu -se conceito 1, a dois cursos ( 12,5 \% ), conceito 2 e a três cursos ( 18,8 \% ), conceito 4.
Nenhum curso ficou Sem Conceito ( SC ) ou recebeu conceito 5 nessa região.
A região Nordeste participou com 36 cursos ou 14,3 \% do total nacional.
Nessa região, 18 cursos ( 50,0 \% em termos regionais ) obtiveram conceito 3, o conceito modal para a região, como já comentado.
Os conceitos 1 e 2 foram atribuídos, respectivamente, a um e a cinco cursos ( 2,8 \% e 13,9 \%, respectivamente ).
Ainda nessa região, oito cursos ( 22,2 \% ) foram avaliados com o conceito 4, e três cursos ( 8,3 \% ), com conceito 5.
Um curso ficou Sem Conceito ( SC ).
Dos 128 cursos participantes da região Sudeste ( 50,8 \% do total nacional ), o conceito 3 foi o conceito modal, como já anteriormente assinalado, tendo sido obtido por 48 cursos ( 37,5 \% ).
O conceito 1 foi atribuído a sete cursos ( 5,5 \% ) e o conceito 2, a 36 cursos ( 28,1 \% ).
Receberam os conceitos 4 e 5, respectivamente, 22 cursos ( 17,2 \% ) e 12 cursos ( 9,4 \% ).
Nessa região, três cursos ficaram Sem Conceito ( SC ).
Os 52 cursos da região Sul corresponderam a 20,6 \% do total nacional.
Como já \textbf{indicado}, houve predominância do conceito 3, atribuído a 42,3 \% dos cursos da região ( conceito modal ), ou seja, atribuído a 22 dos 52 cursos participantes na região Sul.
O conceito 1 foi atribuído a um curso ( 1,9 \% ), o conceito 2, a seis cursos ( 11,5 \% ), o conceito 4, a 17 cursos ( 32,7 \% ) e o conceito 5, a cinco curso ( 9,6 \% ).
Um curso ficou Sem Conceito ( SC ) na região Sul.
Já dos 20 cursos participantes na região Centro-Oeste ( 7,9 \% do total nacional ), oito receberam conceito 3, o conceito modal ( 40,0 \% ), como já fora destacado.
Cinco cursos ( 25,0 \% ), receberam conceito 2.
O conceito 4 foi atribuído a seis cursos ( 30,0 \% ) e um curso ( 5,0 \% ) recebeu conceito 5.
Nessa região, nenhum curso ficou Sem Conceito ( SC ) ou com conceito 1.
O Gráfico 5.1 apresenta a distribuição \textbf{cumulativa} do Conceito Enade segundo Grande Região.
Num gráfico de distribuição \textbf{cumulativa}, em cada valor de \textbf{abscissa} ( eixo x ), o valor apresentado no eixo das \textbf{ordenadas} ( eixo y ) é o valor acumulado da \textbf{variável}.
Por exemplo, a região Norte ( linha azul ) apresenta um pouco mais de 80 \% dos cursos com conceito 3 ou abaixo deste valor.
A região Nordeste, por outro lado, apresenta um pouco menos de 70 \% dos cursos na mesma situação : com conceito 3 ou abaixo.
Na região Norte ( linha azul ), todos os cursos têm conceitos 4 ou abaixo disso, já que a \textbf{poligonal} \textbf{cumulativa} alcança 100 \% para o conceito 4.
Em linhas gerais, num gráfico de distribuição \textbf{cumulativa}, \textbf{poligonais} mais à esquerda representam Grandes Regiões com pior distribuição, e \textbf{poligonais} mais à direita, Grandes Regiões com melhores conceitos.
A região Norte ( linha azul ) apresenta os cursos com conceitos mais baixos, ao passo que a região Sul ( linha \textbf{roxa} ) apresenta as distribuições com valores maiores.
Cursos em IES situadas nas regiões, Nordeste ( linha verde ), Sudeste ( linha vermelha ) e Centro-Oeste ( linha preta ) apresentam uma situação intermediária 5.2 CONCEITOS POR CATEGORIA ADMINISTRATIVA, POR MODALIDADE DE ENSINO E GRANDE REGIÃO A Tabela 5.2 apresenta a distribuição dos cursos de Engenharia Ambiental participantes do Enade/2017, por Categoria Administrativa e Modalidade de ensino, de acordo com os conceitos por eles alcançados, segundo as Grandes Regiões brasileiras.
Dos 252 cursos participantes, 85 ( 33,7 \% ) eram \textbf{ministrados} em Instituições \textbf{Públicas} e 167 ( 66,3 \% ), em Instituições Privadas.
Quanto a Modalidade de ensino, 248 cursos eram da modalidade Presencial e quatro, a Distância.
De acordo com as informações da Tabela 5.2, em termos nacionais, entre Instituições \textbf{Públicas} e Privadas, observa -se que, dos 21 cursos avaliados com conceito 5, 19 eram oferecidos em IES Públicas e dois, em IES Privadas.
Dos 85 cursos participantes de IES Públicas, o conceito 4 foi o valor modal, atribuído a 36 cursos ( 42,4 \% ).
Entre os demais cursos de Instituições \textbf{Públicas} participantes, três cursos ( 3,5 \% ) foram avaliados com conceito 1, quatro cursos ( 4,7 \% ) obtiveram conceito 2, 23 cursos ( 27,1 \% ) receberam conceito 3 e, como já comentado, 19 cursos ( 22,4 \% ) receberam conceito 5.
Nesta categoria, nenhum curso ficou Sem Conceito ( SC ).
Na rede privada, o conceito modal foi 3, atribuído a 81 cursos, o correspondente a 48,5 \% dos 167 cursos da categoria.
Entre os demais cursos participantes, nove ( 5,4 \% ) receberam conceito 1, e 50 ( 29,9 \% ), conceito 2.
O conceito 4 foi atribuído a 20 cursos ( 12,0 \% ) e o conceito 5, a dois cursos ( 1,2 \% ).
Nesta Categoria Administrativa, cinco cursos ( 3,0 \% ) ficaram Sem Conceito ( SC ).
Dos 248 cursos participantes da modalidade presencial, o conceito 3 foi o valor modal, atribuído a 103 cursos ( 41,5 \% ).
Entre os demais cursos nesta modalidade de ensino, 55 cursos ( 22,2 \% ) foram avaliados com conceito 4, 53 cursos ( 21,4 \% ) obtiveram conceito 2, 21 cursos ( 8,5 \% ) receberam conceito 5 e 12 cursos ( 4,8 \% ) receberam conceito 1.
Nesta modalidade quatro cursos ficaram Sem Conceito ( SC ).
Dos quatro cursos a Distância, um recebeu conceito 2, um recebeu conceito 3 e um recebeu conceito 4.
Nessa Modalidade um curso ficou Sem Conceito ( SC ) e nenhum curso recebeu conceito 1 ou 5.
Na análise por região, observa -se que, na região Norte, as Instituições Privadas participaram com 9 cursos ( 56,3 \% do total regional ), dos quais nenhum ficou Sem Conceito ( SC ) e nenhum recebeu conceito 4 ou 5.
O conceito modal para as Instituições Privadas na região foi 3, atribuído a 5 cursos, correspondendo a 55,6 \% dessa categoria na região.
Quatros cursos foram avaliados com conceitos 1 e 3, sendo dois cada um.
As Instituições Públicas participaram com sete cursos na região Norte ( 43,8 \% do total regional ), dos quais seis obtiveram conceito 3 e 4 ( três cursos para cada ), ou seja, o valor modal é contíguo.
O conceito 1 foi atribuído a um curso.
Nenhum curso ficou Sem Conceito ( SC ) ou com conceito 2 ou 5 nessa combinação de categoria e região.
Todos os cursos oferecidos na região Norte eram presenciais.
O conceito modal 3 foi atribuído a oito cursos.
Nenhum curso ficou Sem Conceito ( SC ) ou com conceito 5.
Os demais conceitos foram atribuídos a três cursos ( conceito 1 ), dois cursos ( conceito 2 ) e três cursos ( conceito 4 ).
Na região Nordeste, a rede privada concentrou 20 dos 36 cursos participantes, o equivalente a 55,6 \% do total da região.
Dentre os cursos oferecidos por IES Privadas no Nordeste, 13 cursos foram avaliados com conceito 3, o conceito modal.
Um curso obteve conceito 1, cinco cursos receberam conceito 2 e um curso ficou Sem Conceito ( SC ).
Nessa combinação de categoria e região, nenhum curso alcançou o conceito 4 ou 5.
As Instituições \textbf{Públicas} da região Nordeste participaram com 16 cursos ( 44,4 \% do total da região ).
Desses, oito obtiveram conceito 4, o valor modal, seguidos por outros cinco cursos, que receberam conceito 3.
O conceito 5 foi atribuído a três cursos.
Nenhum curso ficou Sem Conceito ( SC ) ou com conceito 1 ou 2 nessa região.
Todos os cursos na região foram presencias : 36.
Para estes, a moda foi o conceito 3 com 18 cursos.
Os demais cursos presenciais da região receberam os seguintes : conceitos 1 ( um curso ), 2 ( cinco cursos ), 4 ( oito cursos ) e 5 ( três cursos ).
Apenas um curso presencial ficou Sem Conceito ( SC ).
Na região Sudeste, a proporção de cursos da rede privada, 76,6 \%, foi mais elevada do que a registrada nas demais regiões brasileiras, correspondendo a 98 dos 128 cursos participantes.
Nessa categoria e região, o conceito modal foi 3, atribuído a 44 cursos.
Três cursos ficaram Sem Conceito ( SC ).
Os demais foram avaliados com conceito 1 ( cinco cursos ), conceito 2 ( 34 cursos ), conceito 4 ( 11 cursos ) e conceito 5 ( um curso ).
Entre os 30 cursos oferecidos por Instituições \textbf{Públicas} na região Sudeste, o conceito 4 e 5 foi atribuído a 11 cursos, cada, sendo esta a moda contígua nesta categoria e região.
Os demais cursos receberam os conceitos 1 e 2 ( ambos com dois cursos ) e 3 ( quatro cursos ).
Na região Sudeste, nesta categoria, nenhum curso ficou Sem Conceito ( SC ).
Também na região Sudeste, a maioria dos cursos foram presencias : 124.
Para estes, a moda foi o conceito 3 com 47 cursos.
Os demais cursos presenciais da região foram alocados aos conceitos 1 ( 7 cursos ), 2 ( 35 cursos ), 4 ( 21 cursos ) e 5 ( 12 cursos ).
Dois cursos ficaram Sem Conceito ( SC ).
Para os quatro cursos a Distância, um ficou Sem Conceito ( SC ).
Os outros três cursos foram igualmente distribuídos entre os conceitos 2, 3 e 4, com um curso para cada.
Nenhum curso nesta modalidade obteve conceito 1 e 5.
As Instituições Privadas concentraram 31 dos 52 cursos participantes da região Sul, 59,6 \% do total regional, a maior das proporções desta categoria entre as regiões.
Desses, 13 obtiveram conceito 3, o conceito modal.
Nesta combinação de Categoria Administrativa e Grande Região, os conceitos 2 e 4 foram atribuídos a seis e a nove cursos, respectivamente.
Um curso recebeu conceito 5, outro conceito 1 e apenas um curso ficou Sem Conceito ( SC ).
As Instituições \textbf{Públicas} na região Sul participaram com 21 cursos, 40,4 \%.
O conceito modal foi 3, atribuído a nove cursos.
Aos demais cursos foram atribuídos os conceitos 4 ( oito cursos ) e 5 ( quatro cursos ).
Nenhum curso da região Sul oferecido por IES Pública ficou Sem Conceito ( SC ) ou recebeu os conceitos 1 e 2.
Também na região Sul, todos os cursos foram presencias, correspondendo a 52.
Para os cursos presenciais, a moda foi o conceito 3 com 22 cursos.
Os demais cursos presenciais da região foram alocados aos conceitos 1 ( um curso ), 2 ( seis cursos ), 4 ( 17 cursos ) e 5 ( cinco cursos ).
Apenas um curso ficou Sem Conceito ( SC ).
Na região Centro-Oeste, nove dos 20 cursos participantes eram de Instituições Privadas ( 45,0 \% em termos regionais ).
Desses, seis receberam conceito 3, o conceito modal.
Entre os demais, todos os 3 cursos \textbf{restantes} receberam conceito 2.
Nenhum curso ficou Sem Conceito ( SC ) ou com conceito 1, 4 e 5.
Quanto aos 11 cursos oferecidos por Instituições \textbf{Públicas} na região Centro-Oeste ( 55,00 \% do total regional ), seis cursos apresentaram conceito 4, conceito modal, e dois ficaram com conceitos 2 e 3 ( dois cursos, cada ) e um curso obteve conceito 5.
Nenhum curso ficou Sem Conceito ( SC ) ou com conceito 1.
Nesta região, todos os cursos foram presenciais, totalizando 20.
Para os cursos presenciais, a moda foi o conceito 3 com oito cursos.
Os demais foram alocados aos conceitos 2 ( cinco cursos ), 4 ( seis cursos ), 5 ( um curso ).
Nenhum curso presencial ficou Sem Conceito ( SC ) ou com conceito 1.
O Gráfico 5.2 apresenta a distribuição do Conceito Enade segundo a Categoria Administrativa da IES.
Os cursos em IES Públicas ( linha azul ) apresentam uma \textbf{poligonal} mais à direita do que os em IES Privadas ( linha verde ) e, consequentemente, uma distribuição de Conceitos Enade com valores maiores. 5.3 CONCEITOS POR ORGANIZAÇÃO ACADÊMICA E GRANDE REGIÃO Na Tabela 5.3, encontra -se a distribuição dos conceitos atribuídos aos cursos participantes do Enade/2017 na Área de Engenharia Ambiental, por Organização Acadêmica, segundo as Grandes Regiões brasileiras.
Dos 252 cursos de Engenharia Ambiental participantes, 128 eram oferecidos em Universidades, 47 em Centros Universitários, 68 em Faculdades e nove em CEFET/IFET.
Esta distribuição corresponde a, respectivamente, 50,8 \%, 18,7 \%, 27,0 \% e 3,6 \% dos cursos.
De acordo com os dados apresentados, de todos os 21 cursos avaliados com conceito 5, 19 eram vinculados a Universidades.
Os cursos vinculados a esse tipo de Organização Acadêmica tiveram o conceito 3 como conceito modal, atribuído a 44 cursos ( 34,4 \% ).
Os demais cursos vinculados a Universidades receberam os conceitos 1 ( três cursos ), 2 ( 18 cursos ), 4 ( 41 cursos ) e conceito 5 ( 19 cursos, como já mencionado ).
Ficaram Sem Conceito ( SC ), três cursos.
Entre os cursos em Centros Universitários, o conceito modal foi também 3, atribuído a 27 cursos ( 57,5 \% ).
Nenhum curso vinculado a esse tipo de Organização Acadêmica ficou Sem Conceito ( SC ), e os outros cursos receberam os conceitos 1 ( três cursos ), 2 ( 11 cursos ), 4 ( cinco cursos ) e 5 ( um curso ).
Dos 68 cursos mantidos por Faculdades, 30 ( 44,1 \% ) receberam conceito 3, o conceito modal.
O segundo conceito mais frequentemente atribuído aos cursos oferecidos por essa categoria de IES foi o conceito 2, recebido por 25 cursos.
Na sequência, 6 cursos obtiveram conceito 1, quatro cursos obtiveram conceito 4 e um curso, conceito 5.
Dois cursos ficaram Sem Conceito ( SC ).
Dos nove cursos mantidos por CEFET/IFET, seis ( 66,67 \% ) receberam conceito 4, o conceito modal para a Organização.
Três cursos receberam conceito 3, o segundo conceito mais frequentemente atribuído aos cursos oferecidos por essa categoria de IES.
Nenhum curso ficou Sem Conceito ( SC ) ou recebeu o conceito 1, 2 ou 5.
Considerando -se separadamente as regiões brasileiras, nove dos 16 cursos da região Norte são vinculados à Universidades.
Quatro cursos receberam conceito 3, conceito modal.
Os demais cursos receberam os conceitos 1 ( dois cursos ) e 4 ( três cursos ).
Nenhum curso ficou Sem Conceito ( SC ) ou recebeu conceito 2 ou 5 nesta combinação de categoria e região.
Os Centros Universitários da região Norte foram representados por dois cursos, aos quais foram atribuídos os conceitos 2 e 3, para um curso, cada, sendo, portanto, o valor \textbf{bimodal} da Região.
Nenhum curso ficou Sem Conceito ( SC ) ou recebeu os demais conceitos nesta combinação de categoria e região.
As Faculdades participaram com cinco cursos na região Norte que receberam os conceitos 1 ( um curso ), 2 ( um curso ) e 3 ( três cursos, conceito modal ).
Nenhum curso ficou Sem Conceito ( SC ).
Nenhuma Faculdade da região obteve conceito 4 ou 5.
A categoria CEFET/IFET não participou.
Na região Nordeste, as Universidades participaram com 17 dos 36 cursos da Área de Engenharia Ambiental oferecidos na região, dos quais nenhum ficou Sem Conceito ( SC ).
O conceito modal da região foi 3, atribuído a sete cursos.
Receberam os conceitos 4 e 5, respectivamente, seis e três cursos.
Nenhum curso oferecido por Universidades no Nordeste recebeu conceito 1 e, apenas, um curso recebeu conceito 2.
Os Centros Universitários \textbf{contaram} com quatro cursos participante na região Nordeste.
Todos os cursos deste tipo de organização receberam o conceito 3, sendo este o conceito modal.
Nenhum curso oferecido por Centros Universitários no Nordeste ficou Sem Conceito ( SC ) ou recebeu os demais conceitos.
As Faculdades foram representadas por 11 cursos na região Nordeste, em que o valor modal é o conceito 3, atribuído a cinco cursos.
O conceito 1 foi atribuído a um curso, e o conceito 2, a quatro cursos.
Na região Nordeste, nenhum curso oferecido por Faculdades recebeu conceito 4 ou 5 e um curso ficou Sem Conceito ( SC ).
Os CEFET/IFET participaram com quatro cursos, os quais receberam conceitos 3 ( dois cursos ) e 4 ( dois cursos ) sendo, portanto, o valor \textbf{bimodal} dessa organização.
Nenhum curso ficou Sem Conceito ( SC ) nem recebeu conceito 1, 2 ou 5.
Na região Sudeste, as Universidades concentraram 53 dos 128 cursos de Engenharia Ambiental da região.
Entre os cursos em Universidades, nessa região, o conceito modal foi 3, atribuído a 14 cursos.
Os demais cursos receberam os conceitos 1 ( um curso ), 2 ( 12 cursos ), 4 ( 13 cursos ) e 5 ( 11 cursos ).
Dois cursos ficaram Sem Conceito ( SC ).
Os Centros Universitários participaram com 31 cursos na região Sudeste, dos quais 18 obtiveram o conceito modal 3.
O conceito 1 foi atribuído a dois cursos, o conceito 2, a oito cursos e o conceito 4, a três cursos.
Nenhum curso ficou Sem Conceito ( SC ) ou com conceito 5 nessa combinação de organização acadêmica e região.
Já as Faculdades foram representadas por 41 cursos na região Sudeste, dos quais um curso ficou Sem Conceito ( SC ).
Um curso recebeu conceito 5 e, aos demais cursos foram atribuídos os conceitos 1 ( quatro cursos ) e 4 ( três cursos ).
Os conceitos 2 e 3 foram atribuídos a, respectivamente, 16 cursos, sendo o valor \textbf{bimodal} da organização.
Os CEFET/IFET da região Sudeste participaram com três cursos, os quais todos receberam o mesmo conceito, 4.
Nenhum curso ficou Sem Conceito ( SC ), nem recebeu conceitos 1, 2, 3 ou 5.
Dos 52 cursos da região Sul, 36 eram de Universidades, para os quais o conceito modal foi 3, atribuído a 15 cursos.
Nesse tipo de organização, um curso ficou Sem Conceito ( SC ) e nenhum curso recebeu conceito 1.
Os demais cursos receberam os conceitos 2 ( dois cursos ), 4 ( 14 cursos ) e 5 ( quatro cursos ).
Dos nove cursos ligados a Centros Universitários na região Sul, três receberam o o conceito modal 3.
Os demais cursos foram distribuídos da seguinte forma : conceitos 1 e 5 ( um curso cada ) e conceitos 2 e 4 ( dois cursos cada ).
Nenhum curso ligado aos Centros Universitários na região Sul ficou Sem Conceito ( SC ).
Quanto aos sete cursos vinculados a Faculdades na região Sul, quatro cursos obtiveram conceito 3, sendo este o valor modal deste tipo de organização.
Nenhum curso ficou Sem Conceito ( SC ) ou recebeu conceitos 1 ou 5.
Para os conceitos 2 e 4 foram atribuídos, respectivamente, dois e um cursos.
Os CEFET/IFET não tiveram participação na região Sul.
Na região Centro-Oeste, 13 dos 20 cursos eram de Universidades.
Nesse tipo de organização, nenhum curso ficou Sem Conceito ( SC ) ou com conceito 1.
O valor modal foi 4, atribuído a 5 cursos, os demais cursos obtiveram conceitos 2 ( três cursos ) e 3 ( quatro cursos ) e um curso alcançou o conceito máximo, 5.
Os Centros Universitários da região Centro-Oeste \textbf{contaram} com apenas um curso, que obteve o conceito 3.
Dos quatro cursos oferecidos em Faculdades na região Centro-Oeste, apenas dois conceitos foram registrados, 2 e 3, atribuídos a dois cursos, cada, sendo o valor \textbf{bimodal}.
Não houve cursos Sem Conceito ( SC ) ou atribuídos aos demais conceitos nesta organização.
Para CEFET/IFET, dois cursos participaram, sendo atribuídos aos conceitos 3 e 4, um curso cada.
Nesta combinação de região e organização, não houve cursos Sem Conceito ( SC ).
O Gráfico 5.3a apresenta a distribuição do Conceito Enade do curso de Engenharia Ambiental segundo a Organização Acadêmica da IES.
Os cursos em CEFET/IFET ( linha \textbf{roxa} ) apresentam uma \textbf{poligonal} mais à direita do que as demais, \textbf{denotando} uma melhor distribuição de Conceitos Enade.
No outro extremo, temos os cursos oferecidos em Faculdades ( linha vermelha ) com a \textbf{poligonal} mais à esquerda e, consequentemente, uma distribuição de Conceitos Enade com menores valores.
IES de Centros Universitários ( linha verde ) e em Universidades ( linha azul ) apresentam uma situação intermediária, posto que em Universidades ( linha azul ) os cursos apresentam conceitos maiores se comparado com os Centros Universitários ( linha verde ).

% matched lemmas: abscissa, bimodal, contar, cumulativo, denotar, equivaler, indicado, intermédio, ministrar, ordenado, panorama, poligonal, pública, restante, roxa, tabular, variável
\end{textsample}
